\SourceLeaf{64}
zu erhalten. Dann muss $n$ ein echter Bruch sein. Andernfalls
divergieren alle Reihen und sind völlig unbrauchbar.

9. Bisher haben wir uns auf die Reihenentwicklung von $(1+n\cos\varphi)^{\nu}$ beschränkt und die Integration
des vorgelegten Ausdrucks $\mathrm{d}\varphi\,(1+n\cos\varphi)^{\nu}$ beiseitegelassen. Sobald die Reihe selbst gefunden ist,
bereitet die Integration keine Mühe. Denn wenn
\[
(1+n\cos\varphi)^{\nu} = A + A^{(1)}\cos\varphi + A^{(2)}\cos 2\varphi + .\,.\,.\,. + A^{(r)}\cos r\varphi + .\,.\,.\,.
\]
so ergibt sich
\[
\int \mathrm{d}\varphi\,(1+n\cos\varphi)^{\nu} = A\varphi + A^{(1)}\sin\varphi + \tfrac{1}{2}A^{(2)}\sin 2\varphi + .\,.\,. + \frac{1}{r}A^{(r)}\sin r\varphi + .\,.\,.\,.,
\]
Die Koeffizienten dieser Reihe sind aus dem Vorhergehenden bekannt.

10. \textsc{Euler} fügte am Ende von Kapitel VI die Integration des Ausdrucks
\[
\mathrm{d}\varphi\,\log(1+n\cos\varphi),
\]
hinzu. Wie beim vorigen Ausdruck geht er so vor, dass er zunächst mit einem besonderen Kunstgriff $\log(1+n\cos\varphi)$ in eine Reihe
nach Kosinus von Winkelvielfachen entwickelt. Tatsächlich
ist diese Reihe bereits als Spezialfall in derjenigen enthalten, die wir für $(1+n\cos\varphi)^{\nu}$ gefunden haben.
Dazu brauchen wir den bekannten Satz, dass
\[
\log(1+x) = \frac{(1+x)^{0}-1}{0},
\]
und dementsprechend
\[
\log(1+n\cos\varphi) = \frac{(1+n\cos\varphi)^{0}-1}{0}.
\]
\par\noindent\hspace*{\parindent}Im allgemeinen Ausdruck für $(1+n\cos\varphi)^{\nu}$ ziehen wir also vom Anfangskoeffizienten die Einheit ab.
Dann setzen wir in ihm und in allen übrigen $\nu=0$ und dividieren jeden durch $0$ selbst.

$\alpha$) Für $A$, den Anfangskoeffizienten der Reihe zu $(1+n\cos\varphi)^{\nu}$,
lautet nach Abzug seines Anfangsglieds ($=1$) das allgemeine Glied mit Index $h^{\mathrm{tus}}$
\[
= \frac{1}{2^{2h}} . {}^{2h}\mathfrak{B}^{(h)} . {}^{\nu}\mathfrak{B}^{(2h)}\,n^{2h}.
\]
Setzt man darin $\nu=0$, so ergibt sich
\[
\frac{1}{2^{2h}} . \frac{0.(-1).(-2).\,.(-2h+1)}{1.2.\,.\,.\,. 2h} . {}^{2h}\mathfrak{B}^{(h)}\,n^{2h},
\]

